\documentclass[10pt,twocolumn,letterpaper]{article}

\usepackage{cvpr}
\usepackage{times}
\usepackage{epsfig}
\usepackage{graphicx}
\usepackage{amsmath}
\usepackage{amssymb}
\usepackage{booktabs}
\usepackage{multirow}
\usepackage{alltt}
\usepackage{caption} % Added for captionof

\usepackage[pagebackref=true,breaklinks=true,letterpaper=true,colorlinks,bookmarks=false]{hyperref}

\def\cvprPaperID{****} % *** Enter the CVPR Paper ID here
\def\httilde{\mbox{\tt\raisebox{-.5ex}{\symbol{126}}}}

\title{\textbf{A Detailed Progress Report on a Physics-Guided Transformer Network for Underwater Image Enhancement}}
\author{Zahid\\
Project Repository\\
{\tt\small zahid@example.com}
}
\date{\small March 29, 2026}


\begin{document}

\maketitle
\thispagestyle{empty}

%%%%%%%%% ABSTRACT
\begin{abstract}
   This report presents a detailed account of the ongoing research and implementation of a novel deep learning architecture for single-image underwater enhancement. The core of this work is the PhysGuidedUFormer, a hybrid model that synergizes a Transformer-based encoder-decoder backbone with a physics-guided branch for explicit transmission map estimation. We have developed a complete end-to-end software pipeline, encompassing a versatile data loading and augmentation module, the full model architecture with multi-scale feature modulation, and a composite loss function enabling both supervised and perceptually-driven training. An initial experimental phase has been completed, involving a 100-epoch training run on the EUVP dataset. This report provides an in-depth description of the underlying physical model, the network architecture, the composite loss formulation, and a critical analysis of the preliminary quantitative and qualitative results. We highlight key successes, such as the plausible estimation of transmission maps, alongside significant challenges, including low quantitative fidelity and an identified anomaly in our metric calculation pipeline. Finally, we outline a rigorous and detailed roadmap for future work, focusing on debugging, optimization, and comprehensive benchmarking.
\end{abstract}

%%%%%%%%% BODY TEXT
\section{Introduction}

The underwater environment presents a formidable challenge for visual perception. Images captured underwater are subject to severe degradation due to the fundamental properties of light propagation in water. Wavelength-dependent absorption and scattering cause color distortion, typically a dominant blue or green cast, and a loss of contrast and sharpness, often referred to as haze. This degradation not only compromises the aesthetic quality of images but critically impairs the performance of computer vision algorithms for tasks like object detection, semantic segmentation, and 3D reconstruction.

The simplified underwater image formation model (IFM) is often used to describe this process:
\begin{equation}
I(x) = J(x)T(x) + B(1 - T(x))
\label{eq:ifm}
\end{equation}
where $I(x)$ is the observed degraded image at pixel $x$, $J(x)$ is the clear latent image we wish to recover, $B$ is the global background light, and $T(x)$ is the transmission map. The transmission map, $T(x) = e^{-\beta d(x)}$, depends on the scene depth $d(x)$ and the water's attenuation coefficient $\beta$.

Traditional enhancement methods attempt to invert Equation \ref{eq:ifm} by estimating its physical parameters. While effective in some cases, they often rely on strong priors that fail in complex scenes. In contrast, modern data-driven methods, particularly those using deep neural networks, have demonstrated superior performance by learning complex mappings from degraded to clear images. However, these methods can act as "black boxes," potentially producing physically inconsistent results and sometimes failing to generalize to unseen water types.

This project posits that a hybrid approach, combining the representational power of deep networks with explicit physical guidance, can overcome the limitations of both paradigms. We propose a model that learns to enhance images while simultaneously estimating a key physical parameter—the transmission map—and using it to guide the enhancement process at multiple feature scales. This report documents the significant implementation progress toward realizing this vision.

Our main contributions in this ongoing work are:
\begin{enumerate}
    \item The design and full implementation of the \textbf{PhysGuidedUFormer}, a novel architecture integrating a Transformer-based backbone with a physics-inspired branch.
    \item A \textbf{composite loss function} that combines fidelity-based and perception-based objectives to enable both supervised and quality-driven learning.
    \item A complete, end-to-end training and evaluation pipeline, and a preliminary set of experiments that provide critical insights for future development.
\end{enumerate}

\section{Methodology}

\begin{table*}[t]
\centering
\caption{Sequential Architecture of the Main \textbf{PhysGuidedUFormer} Backbone. (Default base C=32, HxW input)}
\label{tab:arch_main}
\resizebox{\textwidth}{!}{\% 
\begin{tabular}{@{}llll@{}}
\toprule
\textbf{Stage} & \textbf{Layer Type} & \textbf{Details} & \textbf{Output Shape} \\
\midrule
\multicolumn{4}{c}{\textbf{Encoder Path}} \\
\midrule
Input & - & Input Image & $H \times W \times 3$ \\
Enc 1 & PatchEmbed (Conv2d) & k=3, s=2, p=1, 3 $\rightarrow$ C & $H/2 \times W/2 \times C$ \\
& TransformerStage & 1 block, C heads & $H/2 \times W/2 \times C$ \\
& $\phi$-Module Fusion & Modulate with T\_map & $H/2 \times W/2 \times C$ \\
\midrule
Enc 2 & PatchEmbed (Conv2d) & k=3, s=2, p=1, C $\rightarrow$ 2C & $H/4 \times W/4 \times 2C$ \\
& TransformerStage & 1 block, 2C heads & $H/4 \times W/4 \times 2C$ \\
& $\phi$-Module Fusion & Modulate with T\_map & $H/4 \times W/4 \times 2C$ \\
\midrule
Enc 3 & PatchEmbed (Conv2d) & k=3, s=2, p=1, 2C $\rightarrow$ 4C & $H/8 \times W/8 \times 4C$ \\
& TransformerStage & 1 block, 4C heads & $H/8 \times W/8 \times 4C$ \\
& $\phi$-Module Fusion & Modulate with T\_map & $H/8 \times W/8 \times 4C$ \\
\midrule
Enc 4 & PatchEmbed (Conv2d) & k=3, s=2, p=1, 4C $\rightarrow$ 8C & $H/16 \times W/16 \times 8C$ \\
& TransformerStage & 1 block, 8C heads & $H/16 \times W/16 \times 8C$ \\
& $\phi$-Module Fusion & Modulate with T\_map & $H/16 \times W/16 \times 8C$ \\
\midrule
\multicolumn{4}{c}{\textbf{Bottleneck}} \\
\midrule
Bottle & TransformerStage & 1 block, 8C heads & $H/16 \times W/16 \times 8C$ \\
& $\phi$-Module Fusion & Modulate with T\_map & $H/16 \times W/16 \times 8C$ \\
\midrule
\multicolumn{4}{c}{\textbf{Decoder Path}} \\
\midrule
Dec 4 & ConvTranspose2d & k=2, s=2, 8C $\rightarrow$ 4C & $H/8 \times W/8 \times 4C$ \\
& Fuse (Concat + Conv2d) & Concat w/ Enc 3; k=1, 8C $\rightarrow$ 4C & $H/8 \times W/8 \times 4C$ \\
& TransformerStage & 1 block, 4C heads & $H/8 \times W/8 \times 4C$ \\
& $\phi$-Module Fusion & Modulate with T\_map & $H/8 \times W/8 \times 4C$ \\
\midrule
Dec 3 & ConvTranspose2d & k=2, s=2, 4C $\rightarrow$ 2C & $H/4 \times W/4 \times 2C$ \\
& Fuse (Concat + Conv2d) & Concat w/ Enc 2; k=1, 4C $\rightarrow$ 2C & $H/4 \times W/4 \times 2C$ \\
& TransformerStage & 1 block, 2C heads & $H/4 \times W/4 \times 2C$ \\
& $\phi$-Module Fusion & Modulate with T\_map & $H/4 \times W/4 \times 2C$ \\
\midrule
Dec 2 & ConvTranspose2d & k=2, s=2, 2C $\rightarrow$ C & $H/2 \times W/2 \times C$ \\
& Fuse (Concat + Conv2d) & Concat w/ Enc 1; k=1, 2C $\rightarrow$ C & $H/2 \times W/2 \times C$ \\
& TransformerStage & 1 block, C heads & $H/2 \times W/2 \times C$ \\
& $\phi$-Module Fusion & Modulate with T\_map & $H/2 \times W/2 \times C$ \\
\midrule
Dec 1 & Interpolate & Scale factor 2 & $H \times W \times C$ \\
& Output Conv2d & k=1, C $\rightarrow$ 3 & $H \times W \times 3$ \\
& Sigmoid & - & $H \times W \times 3$ \\
\bottomrule
\end{tabular}% 
}
\end{table*}

The implemented system is centered on the PhysGuidedUFormer model and its associated learning strategy. The model's architecture is detailed in Table \ref{tab:arch_main} and Table \ref{tab:arch_phys}.

\subsection{Model Architecture: PhysGuidedUFormer}
The architecture is a symmetric U-Net-style encoder-decoder network. Its novelty lies in the use of Transformer-based blocks for feature processing and the integration of a parallel physics-guided branch.

\subsubsection{Encoder Path}
The encoder path consists of four sequential stages, which progressively reduce spatial resolution while increasing feature dimensionality to build a hierarchy of feature representations. Each stage uses a strided convolution for patch embedding and downsampling, followed by a Transformer stage for feature extraction.

\subsubsection{Physics-Guided Branch}
In parallel to the main encoder, a lightweight convolutional U-Net acts as the Physics Branch. Its sole purpose is to estimate a single-channel transmission map $T \in [0, 1]^{H \times W}$ from the input image $I$. This design allows for efficient estimation of the transmission map, which serves as a proxy for scene depth and water clarity. The architecture of this branch is detailed in Table \ref{tab:arch_phys}.

% Non-floating table for PhysicsBranch
\begin{center}
\captionof{table}{Architecture of the \textbf{PhysicsBranch}.}
\label{tab:arch_phys}
\resizebox{\columnwidth}{!}{\% 
\begin{tabular}{@{}lll@{}}
\toprule
\textbf{Layer} & \textbf{Details} & \textbf{Output Shape} \\
\midrule
Input & - & $H \times W \times 3$ \\
Enc 1 & Conv-Block & $H \times W \times 16$ \\
Downsample & AvgPool2d(2) & $H/2 \times W/2 \times 16$ \\
Enc 2 & Conv-Block & $H/2 \times W/2 \times 32$ \\
Downsample & AvgPool2d(2) & $H/4 \times W/4 \times 32$ \\
Enc 3 (Bottle) & Conv-Block & $H/4 \times W/4 \times 64$ \\
Upsample & Interpolate(x2) & $H/2 \times W/2 \times 64$ \\
Dec 2 & Conv-Block & $H/2 \times W/2 \times 32$ \\
Upsample & Interpolate(x2) & $H \times W \times 32$ \\
Dec 1 & Conv-Block & $H \times W \times 16$ \\
Out Conv & Conv2d(k=1) & $H \times W \times 1$ \\
Sigmoid & - & $H \times W \times 1$ \\
\bottomrule
\end{tabular}% 
}
\end{center}


\subsubsection{Multi-Scale Hybrid Fusion}
The fusion of physical guidance and data-driven features is a cornerstone of this model. At each of the four stages in the main encoder and decoder, the estimated transmission map $T$ is downsampled via bilinear interpolation to match the spatial resolution of the feature map $F$ at that stage. A small, stage-specific convolutional network, the \textbf{\phi}-Module}, then processes this resized transmission map to produce a channel-wise attention vector. This vector is applied to the main feature map via element-wise multiplication:
\begin{equation}
    F' = F \odot \phi(T_{\text{resized}})
\end{equation}
This mechanism allows the network to dynamically re-weight its feature channels at every scale based on the estimated physical properties of the scene, enabling a more informed and physically plausible enhancement process.

\subsubsection{Decoder Path}
The decoder path is symmetric to the encoder. It consists of four stages that progressively increase spatial resolution. Each decoder stage involves upsampling via transposed convolution, concatenation with the corresponding skip connection from the encoder, feature fusion, and refinement with a Transformer stage. The output is again modulated by the transmission map before being passed to the next stage.
\subsection{Composite Loss Function}
The network is trained using a composite loss function designed to balance pixel-level accuracy, structural integrity, and perceptual quality. The total loss is a weighted sum of three distinct components:
\begin{equation}
    \mathcal{L} = w_{l1}\mathcal{L}_{L1} + w_{ssim}\mathcal{L}_{SSIM} + w_{nr}\mathcal{L}_{NR}
\end{equation}
\begin{itemize}
    \item \textbf{Fidelity Losses:} For supervised training with paired data, two fidelity losses are used. The Mean Absolute Error, $\mathcal{L}_{L1}$, enforces pixel-level accuracy. The Structural Similarity Index Measure loss, $\mathcal{L}_{SSIM} = 1 - \text{SSIM}(J, \hat{J})$, ensures that the model preserves high-level structural information.
    \item \textbf{Perceptual Quality Loss:} To enable learning from unpaired data and to further improve perceptual quality, a non-reference term, $\mathcal{L}_{NR}$, is included. This term is constructed from differentiable software proxies for two well-known underwater image quality metrics: UCIQE (measuring chroma, saturation, and luminance contrast) and UIQM (measuring colorfulness and sharpness). The loss is formulated as the negative of the predicted quality score, thereby encouraging the model to generate outputs that maximize these perceptual metrics.
\end{itemize}

\section{Experimental Evaluation}
An initial set of experiments was conducted to validate the implementation and establish a preliminary performance baseline.

\subsection{Datasets and Metrics}
The experiments utilize the EUVP dataset, containing paired and unpaired underwater images. Performance is measured using PSNR, SSIM, and the non-reference UCIQE and UIQM metrics. A second, comprehensive synthetic dataset is available for future work.

\subsection{Implementation Details}
The model was trained for 100 epochs using the Adam optimizer. The initial learning rate was set to $1 \times 10^{-4}$ with a batch size of 4. Training was performed on 256$\times$256 random crops. To improve computational efficiency, Automatic Mixed Precision (AMP) was enabled.

\subsection{Results and Analysis}
An evaluation was performed on the EUVP validation set. The mean and standard deviation of the resulting metrics are presented in Table \ref{tab:results_quant}.

% Non-floating table for results
\begin{center}
\captionof{table}{Preliminary Quantitative Results on EUVP Dataset.}
\label{tab:results_quant}
\begin{tabular}{@{}lcc@{}}
\toprule
Metric & Mean & Std. Dev. \\
\midrule
PSNR & 6.43 dB & 1.42 \\
SSIM & 0.263 & 0.11 \\
UCIQE (proxy) & 0.520 & 0.00 \\
UIQM (proxy) & 0.173 & 0.00 \\
\bottomrule
\end{tabular}
\end{center}

The results are decidedly preliminary. The mean PSNR of 6.43 dB is substantially lower than state-of-the-art benchmarks, indicating that the model has not yet converged to an effective solution. The high standard deviation in both PSNR and SSIM suggests that performance is highly inconsistent across different images in the validation set. Most critically, the zero standard deviation for the UCIQE and UIQM proxy metrics points to a significant anomaly in the metric calculation pipeline, as discussed in Section 5.

\section{Qualitative Analysis}
Qualitative results, in the form of side-by-side comparisons of input, output, and ground truth images, were generated. Visual inspection reveals that the model is learning to perform basic color correction, mitigating the strong blue-green cast typical of underwater scenes. The estimated transmission maps are structurally plausible, generally correlating with scene depth (e.g., assigning lower transmission values to distant objects). However, the enhanced images suffer from a noticeable loss of fine detail and the introduction of visual artifacts, which is consistent with the low quantitative fidelity scores.

\section{Current Limitations and Future Work}
The initial experimental phase has been invaluable in identifying critical areas for improvement.

\subsection{Challenges and Limitations}
\begin{enumerate}
    \item \textbf{Sub-optimal Quantitative Performance:} The primary challenge is the poor performance in terms of PSNR and SSIM. This may stem from several factors, including training instability caused by conflicting loss terms, sub-optimal hyperparameter choices, or a potential mismatch between model capacity and data complexity.
    \item \textbf{Metric Calculation Anomaly:} The logging of identical values for all non-reference metrics is a critical bug. It prevents any meaningful assessment of the model's perceptual quality and undermines the role of the $\mathcal{L}_{NR}$ term in the loss function. The software modules responsible for these calculations must be debugged and validated.
\end{enumerate}

\subsection{Roadmap for Future Work}
A clear, structured plan has been formulated to address these challenges and advance the project.

\subsubsection{Short-Term Goals (1-3 weeks)}
\begin{itemize}
    \item \textbf{Debugging:} The immediate priority is to diagnose and fix the non-reference metric calculation anomaly.
    \item \textbf{Performance Tuning:} A systematic investigation into the low fidelity scores will be conducted. This will involve simplifying the loss function to isolate problematic terms and performing a grid search for key hyperparameters like the learning rate and the weights of the composite loss.
    \item \textbf{Training Stabilization:} Techniques such as learning rate scheduling and gradient clipping will be introduced to improve training stability.
\end{itemize}

\subsubsection{Long-Term Goals (4-8 weeks)}
\begin{itemize}
    \item \textbf{Comprehensive Benchmarking:} Once performance is stabilized, the model will be rigorously evaluated on the synthetic dataset and other public benchmarks (UIEB, LSUI).
    \item \textbf{Baseline Comparisons:} At least two established baseline models (e.g., a standard U-Net and FUnIE-GAN) will be implemented and trained under identical conditions to provide a fair and meaningful comparison.
    \item \textbf{Ablation Studies:} A series of ablation studies will be designed to quantify the contribution of each of our model's components. This will include training variants without the physics-guided branch and without the multi-scale fusion mechanism to empirically validate our core architectural hypotheses.
\end{itemize}

{\small
\bibliographystyle{ieee_fullname}
\begin{thebibliography}{9}
\bibitem{ref:uwcnn}
Li, C., Anwar, S., & Porikli, F. (2020). Underwater scene prior inspired deep underwater image and video enhancement. \textit{Pattern Recognition}.

\bibitem{ref:waternet}
Li, C., Guo, C., Ren, W., Cong, R., Hou, J., Kwong, S., & Tao, D. (2019). An underwater image enhancement benchmark dataset and beyond. \textit{IEEE Transactions on Image Processing}.

\bibitem{ref:funiegan}
Islam, M. J., Xia, Y., & Sattar, J. (2020). Fast underwater image enhancement for improved visual perception. \textit{IEEE Robotics and Automation Letters}.
\end{thebibliography}
}

\end{document}
